% ======================================================================
%  UPPGIFT 6: LAGA JULGRANSBELYSNING
% ======================================================================
\section{6 Julgransbelysning}

\uppgiftsbild{6}{Laga julgransbelysning}
{analysera en algoritm \emph{exakt} och bevisa en matchande undre gräns för problemet.}
{\begin{itemize}\setlength{\itemsep}{1pt}
  \item Kostnadsmåttet: lampbyten, inte tester och inte tid.
  \item Algoritmens värsta fall mot problemets undre gräns.
  \item Exakt antal, $2n-1$, inte bara $\Theta(n)$.
\end{itemize}}

\begin{frame}{Uppgift 6: Laga julgransbelysning}
\begin{infobox}[fontupper=\footnotesize]
Jonas julgransbelysning med $n$ seriellt kopplade lampor lyser inte. Det räcker att en lampa är
trasig för att belysningen inte ska lysa. Han har $n$ nya lampor som säkert fungerar, men vill
inte ersätta några hela lampor med nya. En urskruvning följd av en iskruvning tar tiden 1.
Konstruera en algoritm, analysera värstafallet \emph{exakt} och ge en matchande undre gräns.
\end{infobox}
\begin{columns}[T]
\begin{column}{0.55\textwidth}
\centering
\begin{tikzpicture}[scale=0.85]
  \useasboundingbox (0.35,-0.75) rectangle (7.15,1.25);
  \draw[black!45,line width=1pt] (0.5,0) -- (7.0,0);
  \draw[black!45,line width=1pt] (0.5,0) -- (0.5,-0.45) -- (7.0,-0.45) -- (7.0,0);
  \fill[black!55] (1.12,0.02) rectangle (1.38,0.22);
  \filldraw[fill=black!8,draw=kthnavy,line width=0.7pt] (1.25,0.58) circle (0.34);
  \node[font=\footnotesize,text=kthnavy] at (1.25,0.58) {$o_{1}$};
  \node[font=\tiny,text=black!45] at (1.25,-0.22) {1};
  \fill[black!55] (2.37,0.02) rectangle (2.63,0.22);
  \filldraw[fill=black!8,draw=kthnavy,line width=0.7pt] (2.5,0.58) circle (0.34);
  \node[font=\footnotesize,text=kthnavy] at (2.5,0.58) {$o_{2}$};
  \node[font=\tiny,text=black!45] at (2.5,-0.22) {2};
  \fill[black!55] (3.62,0.02) rectangle (3.88,0.22);
  \filldraw[fill=black!8,draw=kthnavy,line width=0.7pt] (3.75,0.58) circle (0.34);
  \node[font=\footnotesize,text=kthnavy] at (3.75,0.58) {$o_{3}$};
  \node[font=\tiny,text=black!45] at (3.75,-0.22) {3};
  \fill[black!55] (4.87,0.02) rectangle (5.13,0.22);
  \filldraw[fill=black!8,draw=kthnavy,line width=0.7pt] (5.0,0.58) circle (0.34);
  \node[font=\footnotesize,text=kthnavy] at (5.0,0.58) {$o_{4}$};
  \node[font=\tiny,text=black!45] at (5.0,-0.22) {4};
  \fill[black!55] (6.12,0.02) rectangle (6.38,0.22);
  \filldraw[fill=black!8,draw=kthnavy,line width=0.7pt] (6.25,0.58) circle (0.34);
  \node[font=\footnotesize,text=kthnavy] at (6.25,0.58) {$o_{5}$};
  \node[font=\tiny,text=black!45] at (6.25,-0.22) {5};
\end{tikzpicture}

\end{column}
\begin{column}{0.42\textwidth}
\small
\begin{itemize}\setlength{\itemsep}{2pt}
  \item \textbf{Kostar 1:} ett byte (ur + i).
  \item \textbf{Gratis:} att titta om det lyser.
  \item \textbf{Förbjudet:} att ersätta en hel gammal lampa.
\end{itemize}
\end{column}
\end{columns}
\varstrip{Separation}{problemet}{kostnadsmåttet}{att vi räknar byten, inte tester}
\end{frame}

\begin{frame}{Viggos algoritm}
\begin{columns}[T]
\begin{column}{0.56\textwidth}
\footnotesize
\begin{algorithmic}[1]
\State byt alla lampor utom den på plats 1 mot nya
\State lägg de gamla lamporna i en säck
\State $i\gets1$
\While{det finns lampor i säcken}
  \If{belysningen lyser}
    \State $i\gets i+1$
    \State skruva ur den nya lampan på plats $i$
  \Else
    \State skruva ur den gamla trasiga lampan på plats $i$
  \EndIf
  \State ta en lampa ur säcken och skruva i den på plats $i$
\EndWhile
\If{belysningen inte lyser}
  \State byt lampan på plats $i$ mot en ny
\EndIf
\end{algorithmic}
\end{column}
\begin{column}{0.41\textwidth}
\begin{infobox}[fontupper=\footnotesize]
\textbf{Invariant} (i början av varje varv): platserna före $i$ har gamla, hela lampor;
plats $i$ har en gammal lampa; platserna efter $i$ har nya.
\end{infobox}
\begin{ahabox}[fontupper=\footnotesize]
Idén: när alla andra lampor är kända och hela avslöjar belysningen exakt en lampa. Alla
gamla lampor provas en i taget, på samma sätt.
\end{ahabox}
\end{column}
\end{columns}
\end{frame}

\begin{frame}{Viggos algoritm steg för steg ($n=5$, trasiga: $o_2$ och $o_5$)}
\only<1>{%
\begin{columns}[T]
\begin{column}{0.62\textwidth}\centering
\begin{tikzpicture}[scale=0.9]
  \useasboundingbox (0.35,-0.75) rectangle (7.15,1.25);
  \draw[black!45,line width=1pt] (0.5,0) -- (7.0,0);
  \draw[black!45,line width=1pt] (0.5,0) -- (0.5,-0.45) -- (7.0,-0.45) -- (7.0,0);
  \fill[black!55] (1.12,0.02) rectangle (1.38,0.22);
  \filldraw[fill=black!8,draw=kthnavy,line width=0.7pt] (1.25,0.58) circle (0.34);
  \node[font=\footnotesize,text=kthnavy] at (1.25,0.58) {$o_{1}$};
  \node[font=\tiny,text=black!45] at (1.25,-0.22) {1};
  \fill[black!55] (2.37,0.02) rectangle (2.63,0.22);
  \filldraw[fill=black!8,draw=kthnavy,line width=0.7pt] (2.5,0.58) circle (0.34);
  \node[font=\footnotesize,text=kthnavy] at (2.5,0.58) {$o_{2}$};
  \node[font=\tiny,text=black!45] at (2.5,-0.22) {2};
  \fill[black!55] (3.62,0.02) rectangle (3.88,0.22);
  \filldraw[fill=black!8,draw=kthnavy,line width=0.7pt] (3.75,0.58) circle (0.34);
  \node[font=\footnotesize,text=kthnavy] at (3.75,0.58) {$o_{3}$};
  \node[font=\tiny,text=black!45] at (3.75,-0.22) {3};
  \fill[black!55] (4.87,0.02) rectangle (5.13,0.22);
  \filldraw[fill=black!8,draw=kthnavy,line width=0.7pt] (5.0,0.58) circle (0.34);
  \node[font=\footnotesize,text=kthnavy] at (5.0,0.58) {$o_{4}$};
  \node[font=\tiny,text=black!45] at (5.0,-0.22) {4};
  \fill[black!55] (6.12,0.02) rectangle (6.38,0.22);
  \filldraw[fill=black!8,draw=kthnavy,line width=0.7pt] (6.25,0.58) circle (0.34);
  \node[font=\footnotesize,text=kthnavy] at (6.25,0.58) {$o_{5}$};
  \node[font=\tiny,text=black!45] at (6.25,-0.22) {5};
\end{tikzpicture}
\end{column}
\begin{column}{0.35\textwidth}
\begin{tikzpicture}[scale=0.9]
  \useasboundingbox (-0.2,-0.75) rectangle (4.9,1.25);
  \filldraw[fill=sand,draw=bark,line width=0.8pt,rounded corners=6pt] (0,-0.5) -- (1.9,-0.5) -- (1.7,0.75) -- (1.25,0.95) -- (0.65,0.95) -- (0.2,0.75) -- cycle;
  \draw[bark,line width=1pt] (0.6,0.95) -- (1.3,0.95);
  \node[font=\tiny,text=bark] at (0.95,1.13) {säcken};
  \node[font=\scriptsize,text width=1.6cm,align=center] at (0.95,0.2) {\textit{tom}};
  \node[font=\scriptsize,anchor=west] at (2.25,0.65) {trasiga: —};
  \node[font=\small\bfseries,text=kthnavy,anchor=west] at (2.25,0.0) {byten: 0};
\end{tikzpicture}
\end{column}
\end{columns}
\medskip
\begin{tcolorbox}[enhanced,halign=left,colback=kthlight!50,colframe=kthsky,boxrule=0.4pt,arc=2pt,left=4pt,right=4pt,top=1pt,bottom=1pt,fontupper=\small,height=1.05cm,valign=center]
Det lyser inte. Minst en lampa är trasig — men vilka?
\end{tcolorbox}}
\only<2>{%
\begin{columns}[T]
\begin{column}{0.62\textwidth}\centering
\begin{tikzpicture}[scale=0.9]
  \useasboundingbox (0.35,-0.75) rectangle (7.15,1.25);
  \draw[kthorange,line width=1pt] (0.5,0) -- (7.0,0);
  \draw[kthorange,line width=1pt] (0.5,0) -- (0.5,-0.45) -- (7.0,-0.45) -- (7.0,0);
  \fill[lampyellow!45] (1.25,0.62) circle (0.5);
  \fill[black!55] (1.12,0.02) rectangle (1.38,0.22);
  \filldraw[fill=lampyellow,draw=kthnavy,line width=0.7pt] (1.25,0.58) circle (0.34);
  \node[font=\footnotesize,text=kthnavy] at (1.25,0.58) {$o_{1}$};
  \node[font=\tiny,text=black!45] at (1.25,-0.22) {1};
  \fill[lampyellow!45] (2.5,0.62) circle (0.5);
  \fill[black!55] (2.37,0.02) rectangle (2.63,0.22);
  \filldraw[fill=lampyellow,draw=kthgreen,line width=1.3pt] (2.5,0.58) circle (0.34);
  \node[font=\footnotesize,text=kthgreen] at (2.5,0.58) {\textbf{N}};
  \node[font=\tiny,text=black!45] at (2.5,-0.22) {2};
  \fill[lampyellow!45] (3.75,0.62) circle (0.5);
  \fill[black!55] (3.62,0.02) rectangle (3.88,0.22);
  \filldraw[fill=lampyellow,draw=kthgreen,line width=1.3pt] (3.75,0.58) circle (0.34);
  \node[font=\footnotesize,text=kthgreen] at (3.75,0.58) {\textbf{N}};
  \node[font=\tiny,text=black!45] at (3.75,-0.22) {3};
  \fill[lampyellow!45] (5.0,0.62) circle (0.5);
  \fill[black!55] (4.87,0.02) rectangle (5.13,0.22);
  \filldraw[fill=lampyellow,draw=kthgreen,line width=1.3pt] (5.0,0.58) circle (0.34);
  \node[font=\footnotesize,text=kthgreen] at (5.0,0.58) {\textbf{N}};
  \node[font=\tiny,text=black!45] at (5.0,-0.22) {4};
  \fill[lampyellow!45] (6.25,0.62) circle (0.5);
  \fill[black!55] (6.12,0.02) rectangle (6.38,0.22);
  \filldraw[fill=lampyellow,draw=kthgreen,line width=1.3pt] (6.25,0.58) circle (0.34);
  \node[font=\footnotesize,text=kthgreen] at (6.25,0.58) {\textbf{N}};
  \node[font=\tiny,text=black!45] at (6.25,-0.22) {5};
\end{tikzpicture}
\end{column}
\begin{column}{0.35\textwidth}
\begin{tikzpicture}[scale=0.9]
  \useasboundingbox (-0.2,-0.75) rectangle (4.9,1.25);
  \filldraw[fill=sand,draw=bark,line width=0.8pt,rounded corners=6pt] (0,-0.5) -- (1.9,-0.5) -- (1.7,0.75) -- (1.25,0.95) -- (0.65,0.95) -- (0.2,0.75) -- cycle;
  \draw[bark,line width=1pt] (0.6,0.95) -- (1.3,0.95);
  \node[font=\tiny,text=bark] at (0.95,1.13) {säcken};
  \node[font=\scriptsize,text width=1.6cm,align=center] at (0.95,0.2) {$o_2$, $o_3$, $o_4$, $o_5$};
  \node[font=\scriptsize,anchor=west] at (2.25,0.65) {trasiga: —};
  \node[font=\small\bfseries,text=kthnavy,anchor=west] at (2.25,0.0) {byten: 4};
\end{tikzpicture}
\end{column}
\end{columns}
\medskip
\begin{tcolorbox}[enhanced,halign=left,colback=kthlight!50,colframe=kthsky,boxrule=0.4pt,arc=2pt,left=4pt,right=4pt,top=1pt,bottom=1pt,fontupper=\small,height=1.05cm,valign=center]
Byt alla utom plats 1 mot nya: 4 byten. De gamla i säcken. Det lyser! \emph{Fantastiskt.} $o_1$ är hel.
\end{tcolorbox}}
\only<3>{%
\begin{columns}[T]
\begin{column}{0.62\textwidth}\centering
\begin{tikzpicture}[scale=0.9]
  \useasboundingbox (0.35,-0.75) rectangle (7.15,1.25);
  \draw[black!45,line width=1pt] (0.5,0) -- (7.0,0);
  \draw[black!45,line width=1pt] (0.5,0) -- (0.5,-0.45) -- (7.0,-0.45) -- (7.0,0);
  \fill[black!55] (1.12,0.02) rectangle (1.38,0.22);
  \filldraw[fill=black!8,draw=kthnavy,line width=0.7pt] (1.25,0.58) circle (0.34);
  \node[font=\footnotesize,text=kthnavy] at (1.25,0.58) {$o_{1}$};
  \node[font=\tiny,text=black!45] at (1.25,-0.22) {1};
  \fill[black!55] (2.37,0.02) rectangle (2.63,0.22);
  \filldraw[fill=black!8,draw=kthnavy,line width=0.7pt] (2.5,0.58) circle (0.34);
  \node[font=\footnotesize,text=kthnavy] at (2.5,0.58) {$o_{2}$};
  \node[font=\tiny,text=black!45] at (2.5,-0.22) {2};
  \fill[black!55] (3.62,0.02) rectangle (3.88,0.22);
  \filldraw[fill=black!8,draw=kthgreen,line width=1.3pt] (3.75,0.58) circle (0.34);
  \node[font=\footnotesize,text=kthgreen] at (3.75,0.58) {\textbf{N}};
  \node[font=\tiny,text=black!45] at (3.75,-0.22) {3};
  \fill[black!55] (4.87,0.02) rectangle (5.13,0.22);
  \filldraw[fill=black!8,draw=kthgreen,line width=1.3pt] (5.0,0.58) circle (0.34);
  \node[font=\footnotesize,text=kthgreen] at (5.0,0.58) {\textbf{N}};
  \node[font=\tiny,text=black!45] at (5.0,-0.22) {4};
  \fill[black!55] (6.12,0.02) rectangle (6.38,0.22);
  \filldraw[fill=black!8,draw=kthgreen,line width=1.3pt] (6.25,0.58) circle (0.34);
  \node[font=\footnotesize,text=kthgreen] at (6.25,0.58) {\textbf{N}};
  \node[font=\tiny,text=black!45] at (6.25,-0.22) {5};
\end{tikzpicture}
\end{column}
\begin{column}{0.35\textwidth}
\begin{tikzpicture}[scale=0.9]
  \useasboundingbox (-0.2,-0.75) rectangle (4.9,1.25);
  \filldraw[fill=sand,draw=bark,line width=0.8pt,rounded corners=6pt] (0,-0.5) -- (1.9,-0.5) -- (1.7,0.75) -- (1.25,0.95) -- (0.65,0.95) -- (0.2,0.75) -- cycle;
  \draw[bark,line width=1pt] (0.6,0.95) -- (1.3,0.95);
  \node[font=\tiny,text=bark] at (0.95,1.13) {säcken};
  \node[font=\scriptsize,text width=1.6cm,align=center] at (0.95,0.2) {$o_3$, $o_4$, $o_5$};
  \node[font=\scriptsize,anchor=west] at (2.25,0.65) {trasiga: —};
  \node[font=\small\bfseries,text=kthnavy,anchor=west] at (2.25,0.0) {byten: 5};
\end{tikzpicture}
\end{column}
\end{columns}
\medskip
\begin{tcolorbox}[enhanced,halign=left,colback=kthlight!50,colframe=kthsky,boxrule=0.4pt,arc=2pt,left=4pt,right=4pt,top=1pt,bottom=1pt,fontupper=\small,height=1.05cm,valign=center]
Plats 2: ny ut, $o_2$ in. Mörkt. \emph{Ånej.} $o_2$ är trasig.
\end{tcolorbox}}
\only<4>{%
\begin{columns}[T]
\begin{column}{0.62\textwidth}\centering
\begin{tikzpicture}[scale=0.9]
  \useasboundingbox (0.35,-0.75) rectangle (7.15,1.25);
  \draw[kthorange,line width=1pt] (0.5,0) -- (7.0,0);
  \draw[kthorange,line width=1pt] (0.5,0) -- (0.5,-0.45) -- (7.0,-0.45) -- (7.0,0);
  \fill[lampyellow!45] (1.25,0.62) circle (0.5);
  \fill[black!55] (1.12,0.02) rectangle (1.38,0.22);
  \filldraw[fill=lampyellow,draw=kthnavy,line width=0.7pt] (1.25,0.58) circle (0.34);
  \node[font=\footnotesize,text=kthnavy] at (1.25,0.58) {$o_{1}$};
  \node[font=\tiny,text=black!45] at (1.25,-0.22) {1};
  \fill[lampyellow!45] (2.5,0.62) circle (0.5);
  \fill[black!55] (2.37,0.02) rectangle (2.63,0.22);
  \filldraw[fill=lampyellow,draw=kthnavy,line width=0.7pt] (2.5,0.58) circle (0.34);
  \node[font=\footnotesize,text=kthnavy] at (2.5,0.58) {$o_{3}$};
  \node[font=\tiny,text=black!45] at (2.5,-0.22) {2};
  \fill[lampyellow!45] (3.75,0.62) circle (0.5);
  \fill[black!55] (3.62,0.02) rectangle (3.88,0.22);
  \filldraw[fill=lampyellow,draw=kthgreen,line width=1.3pt] (3.75,0.58) circle (0.34);
  \node[font=\footnotesize,text=kthgreen] at (3.75,0.58) {\textbf{N}};
  \node[font=\tiny,text=black!45] at (3.75,-0.22) {3};
  \fill[lampyellow!45] (5.0,0.62) circle (0.5);
  \fill[black!55] (4.87,0.02) rectangle (5.13,0.22);
  \filldraw[fill=lampyellow,draw=kthgreen,line width=1.3pt] (5.0,0.58) circle (0.34);
  \node[font=\footnotesize,text=kthgreen] at (5.0,0.58) {\textbf{N}};
  \node[font=\tiny,text=black!45] at (5.0,-0.22) {4};
  \fill[lampyellow!45] (6.25,0.62) circle (0.5);
  \fill[black!55] (6.12,0.02) rectangle (6.38,0.22);
  \filldraw[fill=lampyellow,draw=kthgreen,line width=1.3pt] (6.25,0.58) circle (0.34);
  \node[font=\footnotesize,text=kthgreen] at (6.25,0.58) {\textbf{N}};
  \node[font=\tiny,text=black!45] at (6.25,-0.22) {5};
\end{tikzpicture}
\end{column}
\begin{column}{0.35\textwidth}
\begin{tikzpicture}[scale=0.9]
  \useasboundingbox (-0.2,-0.75) rectangle (4.9,1.25);
  \filldraw[fill=sand,draw=bark,line width=0.8pt,rounded corners=6pt] (0,-0.5) -- (1.9,-0.5) -- (1.7,0.75) -- (1.25,0.95) -- (0.65,0.95) -- (0.2,0.75) -- cycle;
  \draw[bark,line width=1pt] (0.6,0.95) -- (1.3,0.95);
  \node[font=\tiny,text=bark] at (0.95,1.13) {säcken};
  \node[font=\scriptsize,text width=1.6cm,align=center] at (0.95,0.2) {$o_4$, $o_5$};
  \node[font=\scriptsize,anchor=west] at (2.25,0.65) {trasiga: $o_2$};
  \node[font=\small\bfseries,text=kthnavy,anchor=west] at (2.25,0.0) {byten: 6};
\end{tikzpicture}
\end{column}
\end{columns}
\medskip
\begin{tcolorbox}[enhanced,halign=left,colback=kthlight!50,colframe=kthsky,boxrule=0.4pt,arc=2pt,left=4pt,right=4pt,top=1pt,bottom=1pt,fontupper=\small,height=1.05cm,valign=center]
$o_2$ ut, $o_3$ in på samma plats. Det lyser. \emph{Otroligt.} $o_3$ är hel.
\end{tcolorbox}}
\only<5>{%
\begin{columns}[T]
\begin{column}{0.62\textwidth}\centering
\begin{tikzpicture}[scale=0.9]
  \useasboundingbox (0.35,-0.75) rectangle (7.15,1.25);
  \draw[kthorange,line width=1pt] (0.5,0) -- (7.0,0);
  \draw[kthorange,line width=1pt] (0.5,0) -- (0.5,-0.45) -- (7.0,-0.45) -- (7.0,0);
  \fill[lampyellow!45] (1.25,0.62) circle (0.5);
  \fill[black!55] (1.12,0.02) rectangle (1.38,0.22);
  \filldraw[fill=lampyellow,draw=kthnavy,line width=0.7pt] (1.25,0.58) circle (0.34);
  \node[font=\footnotesize,text=kthnavy] at (1.25,0.58) {$o_{1}$};
  \node[font=\tiny,text=black!45] at (1.25,-0.22) {1};
  \fill[lampyellow!45] (2.5,0.62) circle (0.5);
  \fill[black!55] (2.37,0.02) rectangle (2.63,0.22);
  \filldraw[fill=lampyellow,draw=kthnavy,line width=0.7pt] (2.5,0.58) circle (0.34);
  \node[font=\footnotesize,text=kthnavy] at (2.5,0.58) {$o_{3}$};
  \node[font=\tiny,text=black!45] at (2.5,-0.22) {2};
  \fill[lampyellow!45] (3.75,0.62) circle (0.5);
  \fill[black!55] (3.62,0.02) rectangle (3.88,0.22);
  \filldraw[fill=lampyellow,draw=kthnavy,line width=0.7pt] (3.75,0.58) circle (0.34);
  \node[font=\footnotesize,text=kthnavy] at (3.75,0.58) {$o_{4}$};
  \node[font=\tiny,text=black!45] at (3.75,-0.22) {3};
  \fill[lampyellow!45] (5.0,0.62) circle (0.5);
  \fill[black!55] (4.87,0.02) rectangle (5.13,0.22);
  \filldraw[fill=lampyellow,draw=kthgreen,line width=1.3pt] (5.0,0.58) circle (0.34);
  \node[font=\footnotesize,text=kthgreen] at (5.0,0.58) {\textbf{N}};
  \node[font=\tiny,text=black!45] at (5.0,-0.22) {4};
  \fill[lampyellow!45] (6.25,0.62) circle (0.5);
  \fill[black!55] (6.12,0.02) rectangle (6.38,0.22);
  \filldraw[fill=lampyellow,draw=kthgreen,line width=1.3pt] (6.25,0.58) circle (0.34);
  \node[font=\footnotesize,text=kthgreen] at (6.25,0.58) {\textbf{N}};
  \node[font=\tiny,text=black!45] at (6.25,-0.22) {5};
\end{tikzpicture}
\end{column}
\begin{column}{0.35\textwidth}
\begin{tikzpicture}[scale=0.9]
  \useasboundingbox (-0.2,-0.75) rectangle (4.9,1.25);
  \filldraw[fill=sand,draw=bark,line width=0.8pt,rounded corners=6pt] (0,-0.5) -- (1.9,-0.5) -- (1.7,0.75) -- (1.25,0.95) -- (0.65,0.95) -- (0.2,0.75) -- cycle;
  \draw[bark,line width=1pt] (0.6,0.95) -- (1.3,0.95);
  \node[font=\tiny,text=bark] at (0.95,1.13) {säcken};
  \node[font=\scriptsize,text width=1.6cm,align=center] at (0.95,0.2) {$o_5$};
  \node[font=\scriptsize,anchor=west] at (2.25,0.65) {trasiga: $o_2$};
  \node[font=\small\bfseries,text=kthnavy,anchor=west] at (2.25,0.0) {byten: 7};
\end{tikzpicture}
\end{column}
\end{columns}
\medskip
\begin{tcolorbox}[enhanced,halign=left,colback=kthlight!50,colframe=kthsky,boxrule=0.4pt,arc=2pt,left=4pt,right=4pt,top=1pt,bottom=1pt,fontupper=\small,height=1.05cm,valign=center]
Plats 3: ny ut, $o_4$ in. Det lyser. \emph{Häpnadsväckande.}
\end{tcolorbox}}
\only<6>{%
\begin{columns}[T]
\begin{column}{0.62\textwidth}\centering
\begin{tikzpicture}[scale=0.9]
  \useasboundingbox (0.35,-0.75) rectangle (7.15,1.25);
  \draw[black!45,line width=1pt] (0.5,0) -- (7.0,0);
  \draw[black!45,line width=1pt] (0.5,0) -- (0.5,-0.45) -- (7.0,-0.45) -- (7.0,0);
  \fill[black!55] (1.12,0.02) rectangle (1.38,0.22);
  \filldraw[fill=black!8,draw=kthnavy,line width=0.7pt] (1.25,0.58) circle (0.34);
  \node[font=\footnotesize,text=kthnavy] at (1.25,0.58) {$o_{1}$};
  \node[font=\tiny,text=black!45] at (1.25,-0.22) {1};
  \fill[black!55] (2.37,0.02) rectangle (2.63,0.22);
  \filldraw[fill=black!8,draw=kthnavy,line width=0.7pt] (2.5,0.58) circle (0.34);
  \node[font=\footnotesize,text=kthnavy] at (2.5,0.58) {$o_{3}$};
  \node[font=\tiny,text=black!45] at (2.5,-0.22) {2};
  \fill[black!55] (3.62,0.02) rectangle (3.88,0.22);
  \filldraw[fill=black!8,draw=kthnavy,line width=0.7pt] (3.75,0.58) circle (0.34);
  \node[font=\footnotesize,text=kthnavy] at (3.75,0.58) {$o_{4}$};
  \node[font=\tiny,text=black!45] at (3.75,-0.22) {3};
  \fill[black!55] (4.87,0.02) rectangle (5.13,0.22);
  \filldraw[fill=black!8,draw=kthnavy,line width=0.7pt] (5.0,0.58) circle (0.34);
  \node[font=\footnotesize,text=kthnavy] at (5.0,0.58) {$o_{5}$};
  \node[font=\tiny,text=black!45] at (5.0,-0.22) {4};
  \fill[black!55] (6.12,0.02) rectangle (6.38,0.22);
  \filldraw[fill=black!8,draw=kthgreen,line width=1.3pt] (6.25,0.58) circle (0.34);
  \node[font=\footnotesize,text=kthgreen] at (6.25,0.58) {\textbf{N}};
  \node[font=\tiny,text=black!45] at (6.25,-0.22) {5};
\end{tikzpicture}
\end{column}
\begin{column}{0.35\textwidth}
\begin{tikzpicture}[scale=0.9]
  \useasboundingbox (-0.2,-0.75) rectangle (4.9,1.25);
  \filldraw[fill=sand,draw=bark,line width=0.8pt,rounded corners=6pt] (0,-0.5) -- (1.9,-0.5) -- (1.7,0.75) -- (1.25,0.95) -- (0.65,0.95) -- (0.2,0.75) -- cycle;
  \draw[bark,line width=1pt] (0.6,0.95) -- (1.3,0.95);
  \node[font=\tiny,text=bark] at (0.95,1.13) {säcken};
  \node[font=\scriptsize,text width=1.6cm,align=center] at (0.95,0.2) {\textit{tom}};
  \node[font=\scriptsize,anchor=west] at (2.25,0.65) {trasiga: $o_2$};
  \node[font=\small\bfseries,text=kthnavy,anchor=west] at (2.25,0.0) {byten: 8};
\end{tikzpicture}
\end{column}
\end{columns}
\medskip
\begin{tcolorbox}[enhanced,halign=left,colback=kthlight!50,colframe=kthsky,boxrule=0.4pt,arc=2pt,left=4pt,right=4pt,top=1pt,bottom=1pt,fontupper=\small,height=1.05cm,valign=center]
Plats 4: ny ut, $o_5$ in. Mörkt: $o_5$ är trasig. Säcken är tom.
\end{tcolorbox}}
\only<7>{%
\begin{columns}[T]
\begin{column}{0.62\textwidth}\centering
\begin{tikzpicture}[scale=0.9]
  \useasboundingbox (0.35,-0.75) rectangle (7.15,1.25);
  \draw[kthorange,line width=1pt] (0.5,0) -- (7.0,0);
  \draw[kthorange,line width=1pt] (0.5,0) -- (0.5,-0.45) -- (7.0,-0.45) -- (7.0,0);
  \fill[lampyellow!45] (1.25,0.62) circle (0.5);
  \fill[black!55] (1.12,0.02) rectangle (1.38,0.22);
  \filldraw[fill=lampyellow,draw=kthnavy,line width=0.7pt] (1.25,0.58) circle (0.34);
  \node[font=\footnotesize,text=kthnavy] at (1.25,0.58) {$o_{1}$};
  \node[font=\tiny,text=black!45] at (1.25,-0.22) {1};
  \fill[lampyellow!45] (2.5,0.62) circle (0.5);
  \fill[black!55] (2.37,0.02) rectangle (2.63,0.22);
  \filldraw[fill=lampyellow,draw=kthnavy,line width=0.7pt] (2.5,0.58) circle (0.34);
  \node[font=\footnotesize,text=kthnavy] at (2.5,0.58) {$o_{3}$};
  \node[font=\tiny,text=black!45] at (2.5,-0.22) {2};
  \fill[lampyellow!45] (3.75,0.62) circle (0.5);
  \fill[black!55] (3.62,0.02) rectangle (3.88,0.22);
  \filldraw[fill=lampyellow,draw=kthnavy,line width=0.7pt] (3.75,0.58) circle (0.34);
  \node[font=\footnotesize,text=kthnavy] at (3.75,0.58) {$o_{4}$};
  \node[font=\tiny,text=black!45] at (3.75,-0.22) {3};
  \fill[lampyellow!45] (5.0,0.62) circle (0.5);
  \fill[black!55] (4.87,0.02) rectangle (5.13,0.22);
  \filldraw[fill=lampyellow,draw=kthgreen,line width=1.3pt] (5.0,0.58) circle (0.34);
  \node[font=\footnotesize,text=kthgreen] at (5.0,0.58) {\textbf{N}};
  \node[font=\tiny,text=black!45] at (5.0,-0.22) {4};
  \fill[lampyellow!45] (6.25,0.62) circle (0.5);
  \fill[black!55] (6.12,0.02) rectangle (6.38,0.22);
  \filldraw[fill=lampyellow,draw=kthgreen,line width=1.3pt] (6.25,0.58) circle (0.34);
  \node[font=\footnotesize,text=kthgreen] at (6.25,0.58) {\textbf{N}};
  \node[font=\tiny,text=black!45] at (6.25,-0.22) {5};
\end{tikzpicture}
\end{column}
\begin{column}{0.35\textwidth}
\begin{tikzpicture}[scale=0.9]
  \useasboundingbox (-0.2,-0.75) rectangle (4.9,1.25);
  \filldraw[fill=sand,draw=bark,line width=0.8pt,rounded corners=6pt] (0,-0.5) -- (1.9,-0.5) -- (1.7,0.75) -- (1.25,0.95) -- (0.65,0.95) -- (0.2,0.75) -- cycle;
  \draw[bark,line width=1pt] (0.6,0.95) -- (1.3,0.95);
  \node[font=\tiny,text=bark] at (0.95,1.13) {säcken};
  \node[font=\scriptsize,text width=1.6cm,align=center] at (0.95,0.2) {\textit{tom}};
  \node[font=\scriptsize,anchor=west] at (2.25,0.65) {trasiga: $o_2$, $o_5$};
  \node[font=\small\bfseries,text=kthnavy,anchor=west] at (2.25,0.0) {byten: 9};
\end{tikzpicture}
\end{column}
\end{columns}
\medskip
\begin{tcolorbox}[enhanced,halign=left,colback=kthlight!50,colframe=kthsky,boxrule=0.4pt,arc=2pt,left=4pt,right=4pt,top=1pt,bottom=1pt,fontupper=\small,height=1.05cm,valign=center]
$o_5$ ut, en ny in. \emph{Hilvide.} $9=2n-1$ byten.
\end{tcolorbox}}
\vfill
{\footnotesize\color{black!55}Invariant: platserna före $i$ har hela gamla lampor, plats $i$ en gammal lampa, platserna efter $i$ nya. (Utropen är lånade från Alice.)}
\varstrip{Generalisering}{algoritmen}{vilka lampor som är trasiga}{att varje gammal lampa provas en gång}
\end{frame}


\begin{frame}{Exakt analys: det är sista lampan som avgör}
\begin{columns}[T]
\begin{column}{0.5\textwidth}
\small
\[
  \underbrace{(n-1)}_{\text{byt till nya}}\;+\;\underbrace{(n-1)}_{\text{prova säcken}}
  \;+\;\underbrace{[\,o_n\text{ trasig}\,]}_{\text{sista bytet}}
\]
Säcken ger alltid $n-1$ byten. Sista bytet behövs bara om den sista lampan ur säcken är trasig.
\begin{itemize}\setlength{\itemsep}{2pt}
  \item Värsta fallet: $2n-1$.
  \item Bästa fallet: $2n-2$.
\end{itemize}
\end{column}
\begin{column}{0.47\textwidth}
\centering
\footnotesize
\renewcommand{\arraystretch}{1.1}
\begin{tabular}{@{}lc@{}}
\toprule
Trasiga ($n=5$) & Byten\\
\midrule
$o_2$, $o_5$ (exemplet) & 9\\
bara $o_5$ & \textbf{9}\\
alla fem & 9\\
bara $o_1$ & 8\\
bara $o_2$ & 8\\
$o_1,\dots,o_4$ & 8\\
\bottomrule
\end{tabular}
\end{column}
\end{columns}
\vfill
\begin{missbox}[fontupper=\footnotesize]
”Värsta fallet är när alla lampor är trasiga.” Det är \emph{ett} värsta fall. Det räcker att den
sista lampan i säcken är trasig.
\end{missbox}
\varstrip{Separation}{algoritmen}{vilka lampor som är trasiga}{att bara $o_n$ påverkar kostnaden}
\end{frame}

{\kaninbg
\begin{frame}{\kh Mörkret i början är gratis information}
\small
Belysningen lyste inte från början. Alltså är minst en lampa trasig.
\begin{itemize}\setlength{\itemsep}{3pt}
  \item Har $o_1,\dots,o_{n-1}$ alla visat sig hela vet vi att $o_n$ är trasig — utan att prova.
        Låt den nya lampan sitta kvar.
  \item På just det indatat kostar det $2n-3$ i stället för $2n-1$.
  \item Värsta fallet blir inte bättre: med $o_{n-1}$ och $o_n$ trasiga kostar det fortfarande
        $2n-1$.
\end{itemize}
\vfill
\begin{kaninbox}[fontupper=\footnotesize]
Det kan det inte heller bli. Att optimera bort det ena värsta fallet flyttar det bara till ett annat.
Det är vad en undre gräns betyder i praktiken.
\end{kaninbox}
\varstrip{Kontrast}{problemet}{algoritmen: med eller utan knepet}{att bästa fallet kan förbättras, värsta inte}
\end{frame}}

\begin{frame}{Undre gräns: motståndaren}
\small
Påstående: \emph{varje} algoritm behöver $2n-1$ byten i värsta fallet. Viggos motståndare:
\begin{enumerate}\setlength{\itemsep}{3pt}
  \item En lampa avslöjas bara när alla \emph{andra} lampor i belysningen är kända och hela.
  \item Motståndaren låter belysningen vara mörk tills $n-1$ nya lampor sitter i: den sista
        originallampan är trasig. ($n-1$ byten.)
  \item De $n-1$ bortskruvade lamporna har inte avslöjat något. Var och en måste in igen.
        ($n-1$ byten till.)
  \item Den sist provade är också trasig. (1 byte.)
\end{enumerate}
Summa $2n-1$, med bara två trasiga lampor.
\vfill
\begin{missbox}[fontupper=\footnotesize]
”Undre gräns för algoritmen” (rubrik i en äldre bild). Argumentet gäller varje algoritm: det är en
undre gräns för \emph{problemet}. För bara Viggos algoritm räcker ett indata: $o_n$ trasig.
\end{missbox}
\end{frame}

{\kaninbg
\begin{frame}{\kh Ett stramare bevis: två indata som ser likadana ut}
\small
Kör en godtycklig algoritm på $X$: \emph{alla} lampor trasiga. Jämför med $X_j$: bara $o_j$ hel.
\begin{itemize}\setlength{\itemsep}{3pt}
  \item Observationerna skiljer sig bara i läget $S_j$: $o_j$ sitter i och alla andra platser har
        nya lampor. (Då lyser $X_j$, men inte $X$.)
  \item Slutlägena måste skilja sig: $o_j$ ska sitta kvar i $X_j$ men inte i $X$. Körningen på $X$
        måste alltså passera \emph{alla} $S_1,\dots,S_n$.
  \item Minst $n-1$ byten till det första $S_j$. Minst 1 byte mellan två olika $S_j$: $n-1$ gånger.
        Minst 1 byte till slutläget.
\end{itemize}
\[
  (n-1)+(n-1)+1=2n-1\quad\text{byten för varje algoritm, på indatat }X.
\]
\vfill
\begin{kaninbox}[fontupper=\footnotesize]
Viggos algoritm går $S_1\to S_2\to\dots\to S_n$ med ett byte i taget på $X$. Den kan inte göras bättre.
\end{kaninbox}
\end{frame}}

\begin{frame}{Samma idé i uppgift 5 och 6}
\footnotesize
\renewcommand{\arraystretch}{1.12}
\begin{tabularx}{\textwidth}{@{}lXX@{}}
\toprule
& \textbf{Uppgift 5: sökning} & \textbf{Uppgift 6: lampor}\\
\midrule
Kostnad & jämförelser & byten (testerna är gratis)\\
Observation & $<$, $=$ eller $>$ & lyser eller lyser inte\\
Motståndaren håller & den större halvan av utfallen & mörkret så länge som möjligt\\
Undre gräns & $\lceil\log_2(n+1)\rceil$ & $2n-1$\\
Når gränsen & binärsökning & Viggos algoritm\\
\bottomrule
\end{tabularx}
\smallskip

\small En motståndare väljer den dyraste stigen genom algoritmens beslutsträd.
\vfill
\begin{missbox}[fontupper=\footnotesize]
”Binärsök efter den trasiga lampan: $O(\log n)$.” Det sparar \emph{tester}, som är gratis.
Bytena kostar: minst $2n-1$, oavsett metod.
\end{missbox}
\varstrip{Fusion}{bevisidén}{problem, kostnadsmått och observationer samtidigt}{motståndare = värsta stigen}
\end{frame}

\begin{frame}{Vanliga fel i uppgift 6}
\begin{columns}[T]
\begin{column}{0.49\textwidth}
\begin{missbox}[fontupper=\footnotesize]
”$\Theta(n)$.” Uppgiften vill ha exakt antal: $2n-1$. Ordoklass räcker inte när skillnaden
mellan $2n-1$ och $2n-2$ är poängen.
\end{missbox}
\begin{missbox}[fontupper=\footnotesize]
”Byt alla mot nya: $n$ byten.” Hela gamla lampor får inte ersättas.
\end{missbox}
\end{column}
\begin{column}{0.49\textwidth}
\begin{missbox}[fontupper=\footnotesize]
”Min algoritm tar $2n-1$, alltså är det en undre gräns.” Det är algoritmens värsta fall.
Problemets undre gräns kräver ett argument för alla algoritmer.
\end{missbox}
\begin{missbox}[fontupper=\footnotesize]
”Värsta fallet: alla trasiga.” Ett av flera. För Viggos algoritm räcker $o_n$.
\end{missbox}
\end{column}
\end{columns}
\varstrip{Kontrast}{uppgiften}{analysen, rätt och nästan rätt}{exakt, tillåtet, algoritm eller problem}
\end{frame}
